\chapter{Веса Больцмана, статистические суммы и идеальный газ}

\lectureoverview{Мы выводим равновесное распределение вероятностей, максимизируя энтропию при фиксированных нормировке и средней энергии. Затем статистическая сумма позволяет получить среднюю энергию, энтропию и температуру дифференцированием. Первым полным примером служит идеальный газ; затем следуют гравитационное обобщение и предварительный взгляд на флуктуации энергии.}

\section{Задача об энтропии при ограничениях}

Начнем с большой совокупности из \(\mathcal N\) одинаковых, очень слабо взаимодействующих реплик одной и той же подсистемы. Если \(n_i\) реплик находятся в \(i\)-м состоянии, то
\begin{equation}
 p_i=\frac{n_i}{\mathcal N},
 \qquad
 \sum_i p_i=1.
 \label{eq:sm4-occupation-probability}
\end{equation}
Мультиномиальный подсчет из предыдущей лекции дает энтропию на одну реплику,
\begin{equation}
 S[p]=-\sum_i p_i\log p_i.
 \label{eq:sm4-entropy}
\end{equation}
Полезно рассматривать ее как обычную функцию всех переменных \(p_1,p_2,\ldots\). Именно ее простая аддитивная форма делает вычисление возможным.

Совокупность реплик имеет фиксированную полную энергию \(\mathcal N E\), где \(E\) — средняя энергия на одну реплику. Поэтому второе ограничение имеет вид
\begin{equation}
 \sum_i p_iE_i=E.
 \label{eq:sm4-energy-constraint}
\end{equation}
Уровни \(E_i\) относятся к подсистеме; число \(E\) в правой части есть их среднее, взвешенное вероятностями. Максимизация
\eqref{eq:sm4-entropy}
при ограничениях
\eqref{eq:sm4-occupation-probability} и
\eqref{eq:sm4-energy-constraint} выбирает доли заселения, которые с подавляющим превосходством наиболее вероятны.

Для удобства будем минимизировать \(-S=\sum_i p_i\log p_i\). Введем по одному множителю Лагранжа для каждого ограничения:
\begin{equation}
 \begin{aligned}
 \Phi[p;\alpha,\beta]
 &=\sum_i p_i\log p_i\\
 &\quad+\alpha\left(\sum_i p_i-1\right)\\
 &\quad+\beta\left(\sum_i p_iE_i-E\right).
 \end{aligned}
 \label{eq:sm4-augmented-functional}
\end{equation}
Символы \(\alpha\) и \(\beta\) пока не выражают никаких физических предположений. На этом этапе это просто два числа, необходимые для наложения двух ограничений.

\begin{classroomqa}
\audiencequestion{Является ли дополненная функция также функцией \(\alpha\) и \(\beta\)?}
\lecturerresponse{Да. Варьирование по \(\alpha\) и \(\beta\) восстанавливает два ограничения. Однако в непосредственном вычислении мы держим эти множители фиксированными и варьируем вероятности \(p_i\). Их значения затем определяются нормировкой и предписанной средней энергией.
\lecturetimestamp{00:09:32--00:09:49}}
\end{classroomqa}

\begin{figure}[t]
\centering
\includegraphics[width=0.92\textwidth]{lecture_04_figure_02.png}
\caption{Дополненный функционал энтропии после опускания членов, постоянных при вариации вероятностей.}
\lectureframe{00:10:56}
\end{figure}

При дифференцировании по вероятности члены
\(-\alpha\) и \(-\beta E\) не дают вклада. Именно этот сокращенный функционал виден на доске. Они остаются частью
\eqref{eq:sm4-augmented-functional}; для данной производной они просто стали несущественными.

\section{Распределение Больцмана}

Выберем любую одну вероятность. В лекции выбрано \(p_7\), отчасти потому, что \(1\) выглядит слишком особенным, а \(7\) — счастливое число. От этого выбора ничто не зависит. Все нужное дифференцирование сводится к
\begin{equation}
 \frac{d}{dp_i}\bigl(p_i\log p_i\bigr)=\log p_i+1.
\end{equation}
Поэтому
\begin{equation}
 \frac{\partial\Phi}{\partial p_i}
 =\log p_i+1+\alpha+\beta E_i=0.
 \label{eq:sm4-stationarity}
\end{equation}
Решая относительно \(p_i\),
\begin{align}
 \log p_i&=-(1+\alpha)-\beta E_i,\\
 p_i&=e^{-(1+\alpha)}e^{-\beta E_i}.
\end{align}
Первый множитель не зависит от метки состояния. Обозначим его через \(1/Z\):
\begin{equation}
 p_i=\frac{e^{-\beta E_i}}{Z}.
 \label{eq:sm4-boltzmann-weight}
\end{equation}

Вот первый неожиданный результат. Стандартное дифференцирование с ограничениями привело к экспоненциальному закону. Относительные вероятности подчиняются соотношению
\begin{equation}
 \frac{p_i}{p_j}=e^{-\beta(E_i-E_j)}.
 \label{eq:sm4-relative-weight}
\end{equation}
Множитель \(\beta\) управляет тем, насколько быстро вероятность убывает с ростом энергии. Большие положительные \(\beta\) концентрируют распределение около самых низких уровней; малые положительные \(\beta\) оставляют заметный вес и на более высоких уровнях спектра. Мы уже видим, что изменение \(\beta\) меняет среднюю энергию, хотя его отождествление с обратной температурой еще предстоит доказать.

\section{Нормировка и статистическая сумма}

Нормировка определяет не зависящий от состояния множитель:
\begin{equation}
 1=\sum_i p_i
 =\frac{1}{Z}\sum_i e^{-\beta E_i},
\end{equation}
поэтому
\begin{equation}
 Z(\beta)=\sum_i e^{-\beta E_i}.
 \label{eq:sm4-partition-function}
\end{equation}
Это каноническая статистическая сумма. Сначала она появилась как нормировочная константа, но такое описание сильно преуменьшает ее роль. Как только \(Z(\beta)\) известна, термодинамические величины получаются из ее производных.

Уровни энергии в \eqref{eq:sm4-partition-function} являются исходными данными. Они следуют из механики системы вместе с фиксированными внешними условиями: объемом и формой ящика, магнитными или электрическими полями и любыми другими управляющими параметрами. Изменение управляющего параметра меняет спектр и, следовательно, меняет \(Z\). В настоящем вычислении эти параметры, а значит и набор \(\{E_i\}\), считаются фиксированными.

\begin{classroomqa}
\audiencequestion{Что именно представляют оси на экспоненциальном эскизе? Вертикальная координата — это какое-то одно конкретное \(p_i\)?}
\lecturerresponse{Изображенная высота — это вероятность, приписанная уровню, тогда как горизонтальный порядок можно читать либо как метку состояния
\(i\), либо, что физически содержательнее, как энергию \(E_i\). Каждому состоянию соответствует фиксированная энергия. Большие \(\beta\) заставляют вероятность круто убывать с этой энергией; вблизи \(\beta=0\) веса почти постоянны на доступных уровнях.
\lecturetimestamp{00:19:39--00:20:09}}
\end{classroomqa}

\section{Средняя энергия, скрытая в статистической сумме}

Второе ограничение определяет \(\beta\):
\begin{equation}
 E=\sum_i p_iE_i
 =\frac{1}{Z}\sum_i E_i e^{-\beta E_i}.
 \label{eq:sm4-mean-energy}
\end{equation}
Вместо того чтобы заново вычислять эту взвешенную сумму, продифференцируем
\eqref{eq:sm4-partition-function}:
\begin{equation}
 \frac{dZ}{d\beta}
 =-\sum_i E_i e^{-\beta E_i}.
\end{equation}
Сравнение с \eqref{eq:sm4-mean-energy} дает
\begin{equation}
 E(\beta)
 =-\frac{1}{Z}\frac{dZ}{d\beta}
 =-\frac{d\log Z}{d\beta}.
 \label{eq:sm4-energy-logz}
\end{equation}
Логарифм здесь естественен, потому что его производная дает множитель
\(1/Z\).

\begin{figure}[t]
\centering
\includegraphics[width=0.82\textwidth]{lecture_04_figure_03.png}
\caption{На доске одновременно записаны нормировка, ограничение на среднюю энергию и производная статистической суммы.}
\lectureframe{00:26:24}
\end{figure}

\begin{classroomqa}
\audiencequestion{Нижнее уравнение на доске — это новое уравнение для \(Z\) или оно получается из определения над ним?}
\lecturerresponse{Оно получается дифференцированием определения
\(Z=\sum_i e^{-\beta E_i}\). При взятии производной появляется множитель
\(-E_i\). Деление результата на \(Z\) затем воспроизводит взвешенную сумму, определяющую среднюю энергию.
\lecturetimestamp{00:25:58--00:26:37}}
\end{classroomqa}

\begin{classroomqa}
\audiencequestion{Делает ли
\(dZ/d\beta=-E Z\) функцию \(Z\) собственной функцией операции дифференцирования?}
\lecturerresponse{Не в обычном смысле. В уравнении на собственные значения \(Z\) умножалась бы на постоянную, тогда как здесь
\(E=E(\beta)\). Это уравнение задает среднюю энергию как функцию \(\beta\); это не задача на фиксированное собственное значение.
\lecturetimestamp{00:28:49--00:29:43}}
\end{classroomqa}

Простота \eqref{eq:sm4-energy-logz} обманчива. Спектр может быть сложным, и вычисление \(Z\) может оказаться трудным, но как только эта единственная функция найдена, средняя энергия получается одной производной. Такой мотив повторяется по всей статистической механике: элементарные операции, примененные в правильном порядке, дают результаты, которые было бы трудно напрямую угадать.

\begin{classroomqa}
\audiencequestion{\(E\) — это полная энергия всех реплик или средняя энергия одного ящика?}
\lecturerresponse{Это средняя энергия одного ящика. Совокупность реплик имеет полную энергию \(\mathcal N E\). Позже, когда в качестве системы будет выбран целый ящик, содержащий \(N\) частиц газа, \(E\) будет означать энергию всего этого образца газа.
\lecturetimestamp{00:30:55--00:31:10}}
\end{classroomqa}

\begin{classroomqa}
\audiencequestion{Можно ли брать логарифм от
\(\sum_i e^{-\beta E_i}\) почленно?}
\lecturerresponse{Нет. Логарифм превращает произведение в сумму, но не сумму в сумму логарифмов:
\[
 \log\!\left(\sum_i e^{-\beta E_i}\right)
 \ne \sum_i\log\!\left(e^{-\beta E_i}\right).
\]
Это различие пригодится позже, когда у независимых подсистем статистические суммы перемножаются.
\lecturetimestamp{00:31:10--00:31:58}}
\end{classroomqa}

\section{Что фиксировано и чему позволено меняться?}

В одних и тех же формулах присутствуют несколько уровней переменных, и в обсуждении на занятии они намеренно разграничиваются.

\begin{itemize}
 \item Вероятности \(p_i\) — это исходные неизвестные.
 \item Уровни \(E_i\) фиксированы, как только заданы система и ее управляющие параметры.
 \item Множитель \(\alpha\) поглощен нормировкой
 \(Z=e^{1+\alpha}\).
 \item Для предписанной средней энергии \(E\) ограничение фиксирует одно число
 \(\beta\). Однако в семействе равновесных состояний можно использовать либо
 \(E=E(\beta)\), либо \(\beta=\beta(E)\).
\end{itemize}

\begin{classroomqa}
\audiencequestion{Если множители — числа, почему \(\beta\) и
\(E\) теперь рассматриваются как переменные?}
\lecturerresponse{Для одной задачи с ограничениями это действительно числа. Но мы можем сравнивать много равновесных задач с одним и тем же спектром и разными предписанными средними энергиями. В таком семействе \(\beta\) меняется с
\(E\). Фиксация \(E\) определяет \(\beta\); фиксация \(\beta\) определяет
\(E\). Это два способа читать одно и то же соотношение.
\lecturetimestamp{00:31:58--00:33:27}}
\end{classroomqa}

\begin{classroomqa}
\audiencequestion{Исчезло ли \(\alpha\)?}
\lecturerresponse{Исчезло только его имя. Определение
\(e^{-(1+\alpha)}=1/Z\) поглощает его в \(Z\). Сохранять одновременно и \(\alpha\),
и \(Z\) значило бы дважды учитывать один и тот же параметр нормировки.
\lecturetimestamp{00:33:29--00:33:49}}
\end{classroomqa}

Энергия не обязана быть единственной сохраняющейся величиной. Если ограничены полный импульс, число частиц, момент импульса или какая-либо другая величина, ее среднему значению соответствует собственный множитель. Такое расширение меняет ансамбль, но не метод.

\begin{classroomqa}
\audiencequestion{Не делаем ли мы, фиксируя только среднюю энергию, слишком сильных предположений о системе?}
\lecturerresponse{Это предполагает, что в данной задаче не нужно накладывать условие ни на какую другую сохраняющуюся величину. Если, например, фиксирован и полный импульс, нужно добавить ещё одно ограничение и ещё один множитель Лагранжа. Напротив, объём, форма и внешние поля здесь являются управляющими параметрами: они считаются фиксированными и сами определяют энергетические уровни.
\lecturetimestamp{00:33:51--00:35:50}}
\end{classroomqa}

\begin{classroomqa}
\audiencequestion{Должна ли \(\alpha\) фигурировать как ещё один аргумент энергии, и зависит ли \(Z\) от уровней \(E_i\)?}
\lecturerresponse{\(\alpha\) уже содержится в \(Z\), так что независимым аргументом она не является. Статистическая сумма, разумеется, зависит как от фиксированного спектра, так и от \(\beta\). Изменение управляющего параметра изменяет \(E_i\) и, следовательно, изменяет \(Z\); при нынешних производных по \(\beta\) спектр считается фиксированным.
\lecturetimestamp{00:35:52--00:37:10}}
\end{classroomqa}

\begin{classroomqa}
\audiencequestion{Варьируем ли мы отдельные уровни \(E_i\) или среднюю энергию \(E\)?}
\lecturerresponse{Отдельные уровни остаются фиксированными. Средняя энергия меняется потому, что меняются вероятности. Таким образом, семейство параметризуется \(E\) или, эквивалентно, \(\beta\), тогда как один и тот же список уровней \(\{E_i\}\) используется повсюду.
\lecturetimestamp{00:37:12--00:39:04}}
\end{classroomqa}

На практике температуру часто гораздо легче измерить, чем полную энергию газа. Мы можем поместить термометр; взвешивать релятивистскую разность масс и ссылаться на \(E=mc^2\) было бы корректно, но поразительно бесполезно. Чтобы придать \(\beta\) операциональный смысл, нужно теперь связать его с термодинамическим определением температуры.

\section{Энтропия раньше температуры}

Логически энтропия предшествует температуре. Подставим вероятности Больцмана в \eqref{eq:sm4-entropy}:
\begin{equation}
 \log p_i=-\beta E_i-\log Z.
\end{equation}
Тогда
\begin{align}
 S
 &=-\sum_i p_i\log p_i\\
 &=\beta\sum_i p_iE_i+\log Z\sum_i p_i\\
 &=\beta E+\log Z.
 \label{eq:sm4-entropy-logz}
\end{align}
Сумма нормировки даёт конечный множитель, равный единице. Это также первый намёк на свободную энергию Гельмгольца, поскольку позже \(F_{\mathrm H}=-T\log Z\) в единицах, где \(k_B=1\).

\begin{classroomqa}
\audiencequestion{Что случилось с множителем \(1/Z\) в первом члене?}
\lecturerresponse{Он уже входит в \(p_i\). Сумма \(\sum_i p_iE_i\) включает полную нормированную вероятность и по определению является средней энергией \(E\).
\lecturetimestamp{00:43:19--00:43:26}}
\end{classroomqa}

\begin{classroomqa}
\audiencequestion{Не должен ли в \(S=\beta E+\log Z\) стоять минус?}
\lecturerresponse{Нет. Энтропия начинается с минуса, а \(\log p_i=-\beta E_i-\log Z\) содержит ещё два минуса. В итоге они дают два положительных члена, выписанных в \eqref{eq:sm4-entropy-logz}.
\lecturetimestamp{00:45:53--00:46:03}}
\end{classroomqa}

Температура была термодинамически определена как
\begin{equation}
 dE=T\,dS,
 \qquad
 \frac{1}{T}=\frac{dS}{dE},
 \label{eq:sm4-temperature-definition}
\end{equation}
при всех остальных фиксированных управляющих параметрах. Первая попытка продифференцировать \eqref{eq:sm4-entropy-logz} непосредственно по \(E\) быстро загромождается, потому что \(\beta\) тоже зависит от \(E\). Чище взять полный дифференциал:
\begin{align}
 dS
 &=\beta\,dE+E\,d\beta
   +\frac{d\log Z}{d\beta}\,d\beta\\
 &=\beta\,dE+E\,d\beta-E\,d\beta\\
 &=\beta\,dE.
 \label{eq:sm4-entropy-differential}
\end{align}
Сокращение использует \eqref{eq:sm4-energy-logz}. Сравнение \eqref{eq:sm4-entropy-differential} с \eqref{eq:sm4-temperature-definition} даёт
\begin{equation}
 T=\frac{1}{\beta},
 \qquad
 \beta=\frac{1}{T}.
 \label{eq:sm4-beta-temperature}
\end{equation}

\begin{figure}[t]
\centering
\includegraphics[width=0.82\textwidth]{lecture_04_figure_04.png}
\caption{Два члена с \(d\beta\) сокращаются, и остаётся \(dS=\beta\,dE\).}
\lectureframe{00:50:48}
\end{figure}

\begin{classroomqa}
\audiencequestion{Не предполагает ли это сокращение ошибочно, что \(\log Z\) не зависит от энергии или температуры?}
\lecturerresponse{Нет. При фиксированном спектре имеется один независимый термический параметр. Можно написать \(Z=Z(\beta)\), а затем рассматривать \(\beta=\beta(E)\) вдоль равновесного семейства. Полный дифференциал сохраняет всю зависимость от \(d\beta\), и тождество \(d\log Z/d\beta=-E\) как раз и обеспечивает это сокращение.
\lecturetimestamp{00:52:04--00:52:44}}
\end{classroomqa}

На протяжении большей части вычисления температура измеряется в единицах энергии, так что \(k_B=1\). Возвращаясь к лабораторным единицам,
\begin{equation}
 \beta=\frac{1}{k_B T_{\mathrm K}},
 \label{eq:sm4-beta-laboratory}
\end{equation}
где \(T_{\mathrm K}\) — абсолютная температура в кельвинах. Часто удобно писать \(T_{\mathrm e}=k_BT_{\mathrm K}\) и затем использовать \(\beta=1/T_{\mathrm e}\).

\begin{classroomqa}
\audiencequestion{Для некоторого объёма газа система — это одна молекула или вся совокупность?}
\lecturerresponse{Это зависит от того, как определена система. Здесь системой будет весь ящик с газом, так что \(E\) — это энергия всех его молекул. Для достаточно разреженного идеального газа одну молекулу тоже можно рассматривать как подсистему, поскольку частицы не взаимодействуют; факторизация ниже покажет связь между этими двумя описаниями.
\lecturetimestamp{00:53:32--00:54:49}}
\end{classroomqa}

Теперь можно собрать канонические тождества:
\begin{align}
 p_i&=\frac{e^{-\beta E_i}}{Z},&
 Z(\beta)&=\sum_i e^{-\beta E_i},\label{eq:sm4-canonical-pair}\\
 E&=-\frac{d\log Z}{d\beta},&
 S&=\beta E+\log Z,&
 T&=\frac{1}{\beta}.
 \label{eq:sm4-canonical-identities}
\end{align}

\section{Классический идеальный газ}

Рассмотрим \(N\) тождественных точечных частиц массы \(m\) в ящике объёма \(V\). Газ является идеальным, когда взаимодействиями между частицами можно пренебречь. Это может быть хорошим приближением либо потому, что силы слабы, либо потому, что газ настолько разрежен, что тесные сближения редки.

Классическое микросостояние задаётся всеми координатами и импульсами:
\begin{equation}
 (\bm q_1,\ldots,\bm q_N;\bm p_1,\ldots,\bm p_N).
\end{equation}
Имеется \(3N\) координат положения и \(3N\) координат импульса, поэтому состояние — это точка в фазовом пространстве размерности \(6N\). При отсутствии взаимодействий гамильтониан имеет вид
\begin{equation}
 H(\bm p)=\sum_{a=1}^{N}\frac{\bm p_a^2}{2m}
 =\sum_{\mu=1}^{3N}\frac{p_\mu^2}{2m}.
 \label{eq:sm4-ideal-hamiltonian}
\end{equation}

\begin{classroomqa}
\audiencequestion{Имеется ли \(3N\) частиц, и не приводит ли суммирование по всем компонентам импульса к двойному счёту состояний?}
\lecturerresponse{Частиц здесь \(N\), а компонент импульса — \(3N\). Сумма в \eqref{eq:sm4-ideal-hamiltonian} содержит по одному квадратичному члену для каждой компоненты. Затем интеграл по фазовому пространству интегрирует каждую независимую компоненту ровно один раз, так что из-за этого произведения двойного счёта не возникает.
\lecturetimestamp{01:02:29--01:04:04}}
\end{classroomqa}

Для непрерывного классического фазового пространства сумма по состояниям принимает вид
\begin{equation}
 Z_N
 =\frac{1}{N!\,h^{3N}}
 \int_{V^N}d^{3N}q
 \int_{\mathbb R^{3N}}d^{3N}p\,
 e^{-\beta H(\bm p)}.
 \label{eq:sm4-classical-partition}
\end{equation}
В лекции внимание сосредоточено на интегралах, и то же выражение записывается без ячейки фазового пространства \(h^{3N}\). Возвращение этого множителя делает \(Z_N\) безразмерной величиной и отражает квантовое происхождение классического подсчёта состояний. Поскольку \(h\) не зависит от \(\beta\), на производную, используемую ниже для вычисления энергии, он не влияет.

Интеграл по координатам берётся сразу. Гамильтониан не зависит от положения, и каждая частица может находиться где угодно в ящике:
\begin{equation}
 \int_{V^N}d^{3N}q=V^N.
\end{equation}
Интеграл по импульсам факторизуется в \(3N\) одинаковых одномерных гауссовых интегралов:
\begin{equation}
 Z_N
 =\frac{V^N}{N!\,h^{3N}}
 \left[
  \int_{-\infty}^{\infty}dp\,
  e^{-\beta p^2/(2m)}
 \right]^{3N}.
 \label{eq:sm4-factorized-gas}
\end{equation}

Множитель \(1/N!\) отождествляет конфигурации, различающиеся только перестановкой тождественных частиц. Для средней энергии при фиксированных \(N\) и \(V\) он даёт лишь не зависящую от \(\beta\) константу в \(\log Z_N\), поэтому исчезает под действием \(-\partial/\partial\beta\). Однако далеко не во всяком вопросе им можно пренебречь: вместе с \(h^{3N}\) он необходим для безразмерной статистической суммы, экстенсивной энтропии и согласованного химического потенциала.

Теперь положим
\begin{equation}
 x=\sqrt{\frac{\beta}{2m}}\,p,
 \qquad
 dp=\sqrt{\frac{2m}{\beta}}\,dx.
\end{equation}
Тогда гауссов интеграл принимает вид
\begin{align}
 \int_{-\infty}^{\infty}dp\,e^{-\beta p^2/(2m)}
 &=
 \sqrt{\frac{2m}{\beta}}
 \int_{-\infty}^{\infty}dx\,e^{-x^2}\\
 &=\sqrt{\frac{2\pi m}{\beta}}.
 \label{eq:sm4-gaussian}
\end{align}

\begin{figure}[t]
\centering
\includegraphics[width=0.82\textwidth]{lecture_04_figure_05.png}
\caption{Интеграл по импульсу сводится к стандартному гауссову интегралу одним масштабированием.}
\lectureframe{01:13:12}
\end{figure}

Подстановка в \eqref{eq:sm4-factorized-gas} даёт
\begin{equation}
 Z_N
 =\frac{V^N}{N!\,h^{3N}}
 \left(\frac{2\pi m}{\beta}\right)^{3N/2}.
 \label{eq:sm4-ideal-partition}
\end{equation}

\begin{figure}[t]
\centering
\includegraphics[width=0.78\textwidth]{lecture_04_figure_06.png}
\caption{На доске — окончательный вид статистической суммы идеального газа. Ячейка фазового пространства \(h^{3N}\) восстановлена в \eqref{eq:sm4-ideal-partition}.}
\lectureframe{01:14:38}
\end{figure}

Удобно записать тот же результат через тепловую длину волны де Бройля
\begin{equation}
 \lambda_T=\frac{h}{\sqrt{2\pi mT}}
 \qquad (k_B=1).
\end{equation}
Тогда
\begin{equation}
 Z_N=\frac{1}{N!}\left(\frac{V}{\lambda_T^3}\right)^N.
 \label{eq:sm4-thermal-wavelength-form}
\end{equation}

При макроскопических \(N\) формула Стирлинга даёт
\begin{equation}
 N!\simeq N^Ne^{-N}\sqrt{2\pi N}.
\end{equation}
Если сохранить только ведущую экспоненциальную зависимость и ввести \(\rho=N/V\),
\begin{equation}
 \frac{V^N}{N!}
 \simeq
 \frac{1}{\sqrt{2\pi N}}
 \left(\frac{e}{\rho}\right)^N.
 \label{eq:sm4-density-rewrite}
\end{equation}
Это аккуратная перегруппировка по плотности, сделанная в лекции. Множитель с квадратным корнем даёт подведущий вклад в \(\log Z_N\), но здесь он выписан, а не молча отброшен.

\begin{classroomqa}
\audiencequestion{Если \(1/N!\) не зависит от \(\beta\), значит ли это, что он вообще ни на что не влияет?}
\lecturerresponse{Для текущей производной средней энергии при фиксированных \(N\) и \(V\) он роли не играет, потому что добавляет не зависящую от \(\beta\) константу к \(\log Z_N\). Для абсолютной энтропии и для изменений числа частиц он важен. Не следует смешивать это ограниченное утверждение о производной с более широким утверждением о подсчёте состояний.
\lecturetimestamp{01:27:01--01:27:51}}
\end{classroomqa}

\section{Энергия и равнораспределение}

Произведения неудобно дифференцировать; логарифмы превращают их в суммы. Из \eqref{eq:sm4-ideal-partition} получаем
\begin{equation}
 \log Z_N
 =-\frac{3N}{2}\log\beta+\text{\(\beta\)-independent terms}.
 \label{eq:sm4-log-ideal-partition}
\end{equation}
Следовательно,
\begin{equation}
 E
 =-\frac{\partial\log Z_N}{\partial\beta}
 =\frac{3N}{2\beta}
 =\frac{3N}{2}T.
 \label{eq:sm4-ideal-energy}
\end{equation}
Энергия на частицу равна
\begin{equation}
 \frac{E}{N}=\frac{3}{2}T
 \qquad (k_B=1),
\end{equation}
или, в лабораторных единицах,
\begin{equation}
 \frac{E}{N}=\frac{3}{2}k_BT_{\mathrm K}.
 \label{eq:sm4-ideal-energy-lab}
\end{equation}

\begin{figure}[t]
\centering
\includegraphics[width=0.84\textwidth]{lecture_04_figure_07.png}
\caption{В энергию идеального газа вносит вклад только член \(-3N\log\beta/2\), что даёт \(E/N=3T/2\).}
\lectureframe{01:25:28}
\end{figure}

\begin{classroomqa}
\audiencequestion{Где в формулу для энергии входит постоянная Больцмана?}
\lecturerresponse{У теоретика температура \(T\) имеет размерность энергии. Чтобы пользоваться кельвинами, замените её на \(k_BT_{\mathrm K}\), получая \(E/N=3k_BT_{\mathrm K}/2\). Малое значение \(k_B\) переводит макроскопическое по величине численное значение температуры в микроскопический энергетический масштаб одной молекулы.
\lecturetimestamp{01:24:58--01:26:01}}
\end{classroomqa}

Число 3 возникает из трёх направлений пространства. Множитель 1/2 возникает из корневой зависимости каждого гауссова интеграла от \(1/\beta\). Каждая независимая квадратичная координата импульса даёт вклад \begin{equation}
 \frac{1}{2}T
\end{equation} в среднюю энергию. В \(d\) пространственных измерениях тот же довод даёт \(E/N=dT/2\). Это первое появление теоремы о равнораспределении.

Расчёт энтропии, заданный в конце лекции, следует из \(S=\beta E+\log Z_N\). Используя приближение Стирлинга,
\begin{equation}
 \frac{S}{N}
 \simeq
 \log\!\left[
  \frac{V}{N}
  \left(\frac{2\pi mT}{h^2}\right)^{3/2}
 \right]
 +\frac{5}{2},
 \label{eq:sm4-sackur-tetrode}
\end{equation}
с точностью до подведущей поправки при конечном \(N\). Восстановление \(k_B\) умножает правую часть на \(k_B\) и заменяет \(T\) на \(k_BT_{\mathrm K}\). Результат ясно показывает, почему нормировочные множители, которые выпали из выражения для энергии, необходимы для энтропии.

\section{Газ в гравитационном поле}

Вопрос о том, что будет, если поставить ящик на нейтронную звезду, подводит к гравитационному обобщению. Прежде чем переходить к нейтронной звезде, рассмотрим однородное поле \(g\) в обычной комнате. Выберем \(y=0\) на дне ящика с площадью основания \(A\) и высотой \(L\). Гамильтониан принимает вид
\begin{equation}
 H=\sum_{a=1}^{N}
 \left(\frac{\bm p_a^2}{2m}+mgy_a\right).
 \label{eq:sm4-gravity-hamiltonian}
\end{equation}
Интегралы по горизонтальным координатам по-прежнему дают \(A^N\), тогда как каждая вертикальная координата даёт вклад
\begin{equation}
 I_y(L)
 =\int_0^Ldy\,e^{-\beta mgy}
 =\frac{1-e^{-\beta mgL}}{\beta mg}.
 \label{eq:sm4-finite-height-integral}
\end{equation}

\begin{classroomqa}
\audiencequestion{Должен ли вертикальный интеграл действительно простираться до бесконечности или только до верхней границы ящика?}
\lecturerresponse{Для конечного ящика он должен идти от \(0\) до его высоты \(L\), как в \eqref{eq:sm4-finite-height-integral}. Переход к бесконечному верхнему пределу — это последующее упрощение, моделирующее атмосферу над поверхностью, а не конечную комнату.
\lecturetimestamp{01:30:25--01:31:32}}
\end{classroomqa}

Поэтому полная статистическая сумма при конечной высоте равна
\begin{equation}
 Z_N(L)
 =
 \frac{1}{N!\,h^{3N}}
 \left[
  A\,
  \frac{1-e^{-\beta mgL}}{\beta mg}
  \left(\frac{2\pi m}{\beta}\right)^{3/2}
 \right]^N.
 \label{eq:sm4-gravity-partition-finite}
\end{equation}
Вертикальный множитель зависит от \(\beta\), так что он должен давать вклад в энергию. Действительно,
\begin{equation}
 \langle U\rangle
 =N\left[
  \frac{1}{\beta}
  -\frac{mgL}{e^{\beta mgL}-1}
 \right].
 \label{eq:sm4-finite-gravity-energy}
\end{equation}
Этим завершается небольшой фрагмент интегрирования, который на лекции сначала оставляют как досадное вычисление.

Для атмосферы бесконечной высоты, \(L\to\infty\), и
\begin{equation}
 I_y(\infty)=\frac{1}{\beta mg}.
\end{equation}
Статистическая сумма принимает вид
\begin{equation}
 Z_N^{(\infty)}
 =
 \frac{1}{N!\,h^{3N}}
 \left[
  \frac{A}{\beta mg}
  \left(\frac{2\pi m}{\beta}\right)^{3/2}
 \right]^N.
 \label{eq:sm4-gravity-partition-infinite}
\end{equation}
Её логарифм, зависящий от \(\beta\), равен
\begin{equation}
 \log Z_N^{(\infty)}
 =-\frac{3N}{2}\log\beta-N\log\beta+\mathrm{const}.
\end{equation}
Первый член даёт кинетическую энергию \(3NT/2\); второй даёт среднюю гравитационную потенциальную энергию \(NT\). Следовательно,
\begin{equation}
 E_{\mathrm{tot}}=\frac{5N}{2}T
\end{equation}
для этой идеализированной неограниченной сверху атмосферы.

\begin{figure}[t]
\centering
\includegraphics[width=0.83\textwidth]{lecture_04_figure_08.png}
\caption{Статистическая сумма для бесконечной атмосферы добавляет \(-N\log\beta\) к \(\log Z\), что и является источником средней потенциальной энергии.}
\lectureframe{01:36:38}
\end{figure}

\begin{classroomqa}
\audiencequestion{Потребовало бы очень сильное гравитационное поле ещё одного множителя Лагранжа?}
\lecturerresponse{Нет, не просто потому, что поле сильное. Гравитация изменяет гамильтониан и тем самым энергетические уровни или вес в фазовом пространстве. Остаётся тот же энергетический множитель \(\beta\). Нейтронная звезда вносит дополнительные физические осложнения, но первой модификацией всё равно остаётся включение гравитационной потенциальной энергии в \(H\).
\lecturetimestamp{01:27:58--01:29:24}}
\end{classroomqa}

\section{Три вопроса в заключение}

\begin{classroomqa}
\audiencequestion{Какое уравнение задаёт распределение Больцмана и что оно утверждает о равновесии?}
\lecturerresponse{Распределение имеет вид
\[
 p_i=\frac{e^{-\beta E_i}}{Z}.
\]
Это распределение максимальной энтропии при условиях нормировки и фиксированной средней энергии. Среднюю энергию нужно вычислять по тем же вероятностям, а условие нормировки всё равно должно быть наложено, хотя в окончательной формуле оно выглядит почти автоматическим.
\lecturetimestamp{01:38:14--01:38:45}}
\end{classroomqa}

Эта процедура может казаться механической: записать статистическую сумму, взять логарифм и продифференцировать. И всё же получающиеся уравнения продолжают удивлять. Опыт подсказывает, какая простая операция должна следовать за какой, но от этого сама тема не становится менее занимательной.

\begin{classroomqa}
\audiencequestion{Можно ли применить эти идеи об энтропии к коллапсирующей звезде или чёрной дыре? Как энтропия может расти, если вещество занимает меньшую область?}
\lecturerresponse{Со звёздами трудно главным образом потому, что подробное описание их вещества сложно. Для чёрной дыры имеется более ясный эталон: её энтропия Бекенштейна--Хокинга является предельной гравитационной энтропией при заданных макроскопических ограничениях. Коллапс может увеличивать энтропию, пока не будет достигнуто это значение для чёрной дыры; для её дальнейшего увеличения нужно изменить доступную энергию или размер.
\lecturetimestamp{01:39:29--01:40:36}}
\end{classroomqa}

\begin{classroomqa}
\audiencequestion{Может ли статистическая сумма также вычислять флуктуации энергии?}
\lecturerresponse{Да. Одна производная даёт среднюю энергию, а вторая производная — её дисперсию:
\begin{align*}
 \frac{\partial^2\log Z}{\partial\beta^2}
 &=\langle E^2\rangle-\langle E\rangle^2\\
 &=\operatorname{Var}(E).
\end{align*}
В энергетических единицах это связано с теплоёмкостью соотношением \(\operatorname{Var}(E)=T^2C_V\). В лабораторных единицах — \(\operatorname{Var}(E)=k_BT_{\mathrm K}^2C_V\). Это соотношение для флуктуаций и будет следующим вычислением.
\lecturetimestamp{01:40:37--01:41:53}}
\end{classroomqa}

Флуктуации отличают статистическую механику от чисто непрерывного термодинамического описания. Их систематическое изучение, связанное прежде всего с Гиббсом и Эйнштейном, выявляет конечные микроскопические составляющие, стоящие за гладкими макроскопическими законами. Именно с этого начинается следующая лекция.